\subsection{Nondegeneracy of the maximiser}\label{sec:seven-entropy-nondegenerate}
Having identified the global maximum and its equality orbit, we now quantify the entropy loss near that orbit.
The constrained Hessian at \(\phi\) is negative definite transverse to the symmetry orbit, giving
an explicit quadratic basin.

\begin{theorem}[Strict full-space local entropy maximum; exact certificate SEH-1--SEH-3 with an analytic perturbation bound]
\label{thm:entropy-seven-local}
The state \(\psi_{x_*}\) is a \emph{strict local maximiser} of
\(E_3\) modulo local unitaries and global phase. More precisely,
for every nonzero real unit-sphere tangent vector \(v\) orthogonal to the
local-unitary and phase tangent space at \(\psi_{x_*}\),
\begin{equation}\label{eq:entropy-seven-hessian-gap}
   D^2(E_3|_{\|z\|=1})_{\psi_{x_*}}[v,v]
      <-\frac{239}{1000}\,\|v\|^2.
\end{equation}
The \(22\) phase and local-unitary symmetry infinitesimals have
zero constrained Hessian.
\end{theorem}
\begin{proof}
Extend \(E_3(z)\) off the unit sphere by the same formula
\(-35^{-1}\sum_{|S|=3}\operatorname{tr}(R_S\log_2 R_S)\), with \(R_S\) not renormalised:
for unit \(z\) and \(r>0\), \(E_3(rz)=r^2E_3(z)-r^2\log_2r^2\). Hence \(E_3\) is the mean entropy only on the
unit sphere, and its radial derivative \(\frac12\frac{d}{dr}E_3(rz)|_{r=1}=E_3(z)-1/\ln2\) is the multiplier \(\lambda\) below.
For a three-qubit cut flatten \(z\) to an \(8\times16\) matrix \(M\)
and put \(R=MM^\dagger\). For a variation \(v\) with flattening \(\dot M\),
write \(\dot R=\dot MM^\dagger+M\dot M^\dagger\), and denote the Fréchet
derivative of the matrix logarithm by
\[
 {\cal L}_R(H)=\int_0^\infty(R+tI)^{-1}H(R+tI)^{-1}\,dt.
\]
The unnormalised second variation of the mean entropy is
\begin{equation}\label{eq:entropy-seven-H-form}
 D^2 E_3(z)[v,v]
 =-\frac1{35\ln2}\sum_{|S|=3}
 \left[2\operatorname{tr}\dot M^\dagger(\ln R+I)\dot M
 +\operatorname{tr}\dot R\,{\cal L}_R(\dot R)\right].
\end{equation}
At the stationary unit vector put
\(\lambda=\langle z,\partial_{\bar z} E_3(z)\rangle
=E_3(z)-1/\ln2\).
The constrained Hessian is \(H=D^2E_3-2\lambda I_{256}\);
it splits into a real and imaginary \(128\times128\) block
because \(z\) is real. Its phase/LU symmetry generators, given
by \(iz\) and \(i\sigma_a^{(j)}z\), span \(G\).

For the certificate take the nearby \emph{rational} vector
\(z_0=a_0U+b_0V\), with \(a_0=16934604/10^8\),
\(b_0=26425139/10^8\). Exact integer signs in
\eqref{eq:entropy-seven-star}, evaluated at
\(698287976009/10^{13}\) and \(698287976010/10^{13}\), bracket
\(x_*\). Squaring rational endpoints proves
\(|a-a_0|,|b-b_0|<5\cdot10^{-9}\) and hence
\(\|\psi_{x_*}-z_0\|<3\cdot10^{-8}\) (SEH-1).
Every reduced eigenvalue along this short segment is at least
\(1/12\), and each difference between two \emph{distinct}
reduced eigenvalues at \(z_0\) exceeds \(1/64\); both are
exact rational checks.

At \(z_0\) the spectral projectors are rational: if \(\ell\)
ranges over the two or three distinct eigenvalues of \(R\), then
\[
 P_\ell=\prod_{h\ne\ell}\frac{R-hI}{\ell-h},\qquad
 \ln R=\sum_\ell(\ln\ell)P_\ell,\qquad
 {\cal L}_R(H)=\sum_{\ell,h}
 \frac{\ln\ell-\ln h}{\ell-h}P_\ell H P_h,
\]
with diagonal coefficient \(1/\ell\). SEH-2 uses rational proxies
for the four logarithms and \(1/\ln2\), each with error less than
\(10^{-13}\), bounded by exact rational \(\operatorname{arctanh}\)
series. It forms a rational fixed-point matrix \(\widehat B\)
approximating \(B=-H+4GG^T\), including the symmetry penalty,
with entrywise rounding checked by integer arithmetic.

The exact Hessian at \(z_0=a_0U+b_0V\) respects the symmetry-adapted orthogonal integer decomposition of Theorem~\ref{thm:icolocal}.
Rounding could disrupt this structure in \(\widehat B\), so SEH-2 checks the decomposition exactly for \(\widehat B\). The decomposition consists of the six computational-basis sectors~\eqref{eq:ico-hessian-six-blocks}, the disjoint-pair projectors \(P_2,P_3,P_{-1},P_{-3}\) on \(\mathcal H_{\mathrm{supp}}\) and \(\mathcal H_{\mathrm{flip}}\), and the triple-incidence projectors \(F_1,F_2,P_{8,40}\) on \(\mathcal H_{\mathrm{odd}}\). In both \(128\times128\) blocks, SEH-2 verifies by exact integer multiplication that the \(16\) isotypic blocks (sizes \(1\times1\), \(2\times2\), and \(4\times4\), with diagonal integer Gram matrices \(D_{\mathrm{blk}}\)) span \(\mathbb R^{128}\) and are exact invariant subspaces of \(\widehat B\), and that every \(1\times1\), \(2\times2\), and \(4\times4\) symmetric integer matrix \(V_{\mathrm{blk}}^{\mathsf T}(\widehat B-0.239395\,I_{128})V_{\mathrm{blk}}\) has strictly positive leading principal minors. Hence \(\lambda_{\min}(\widehat B)>0.239395\) with no floating-point basis or Gram conditioning step.

It remains to bound \(B(\psi_{x_*})-\widehat B\) analytically. Write \(\varepsilon=\|z_0-\psi_{x_*}\|<3\cdot10^{-8}\).
Using \(R\succeq I/12\) along the segment,
\[
 \|R(z_0)-R(z_*)\|_{\rm op}\le3\varepsilon,\quad
 \|\ln R(z_0)-\ln R(z_*)\|_{\rm op}\le36\varepsilon,\quad
 \|{\cal L}_{R(z_0)}-{\cal L}_{R(z_*)}\|_{F\to F}
 \le432\varepsilon.
\]
Also \(\|\dot R(v)\|_F\le3\|v\|\) and
\(\|\dot R_{z_0}(v)-\dot R_{z_*}(v)\|_F
\le2\varepsilon\|v\|\). Inserting these estimates in
\eqref{eq:entropy-seven-H-form} with \(\|\dot M\|_F=\|v\|=1\) and
\(\|{\cal L}_R\|\le12\), each cut changes by at most
\(2\cdot36\varepsilon+2\cdot(2\varepsilon)\cdot12\cdot3+3\cdot432\varepsilon\cdot3
=4104\varepsilon\), hence the mean by at most
\(4104\varepsilon/\ln2<5.93\cdot10^3\varepsilon\). Parametrise the segment from $z_*$ to $z_0$ over $[0,1]$;
then $\|\dot R\|_F\le3\varepsilon$ along it. Hence
\(|\operatorname{tr}\dot R(\ln R+I)|\le\|\dot R\|_F\|\ln R+I\|_F<3\sqrt8\cdot\frac32\,\varepsilon\) per cut and
\(|\|z_0\|^2-1|<3\varepsilon\), so \(\lambda(z)=E_3(z)-\|z\|^2/\ln2\) obeys \(2|\lambda(z_0)-\lambda(z_*)|<50\varepsilon\); the \(22\) columns \(iz,i\sigma_a^{(j)}z\)
of \(G\) have norm \(\|z\|\), so \(4\|G(z_0)G(z_0)^T-G(z_*)G(z_*)^T\|_{\rm op}\le8\cdot22(1+\varepsilon)\varepsilon<180\varepsilon\).
These add less than \(4\cdot10^2\varepsilon\). The operator-norm change is therefore below
\(10^4\varepsilon\). This budget is analytic; SEH-3 checks only the resulting margin.
The four rational logarithm enclosures, their divided differences,
and the explicit \(10^{-10}\)-grid rounding used to assemble
\(\widehat B\) add less than \(10^{-5}\) in operator norm.
Consequently
\[
 \|B(\psi_{x_*})-\widehat B\|_{\rm op}
 <10^4(3\cdot10^{-8})+10^{-5}=0.00031.
\]
The exact block positivity bound now implies \(B(\psi_{x_*})\succ0\), with
\(\lambda_{\min}(B)>0.239395-0.00031=0.239085>0.239\) (SEH-3).
By Proposition~\ref{prop:entropy-seven-stationarity} and LU/phase
invariance, \(H\) annihilates every vector in \(G\). On \(G^\perp\)
the penalty vanishes, so \(-H=B\succ0\) and
\eqref{eq:entropy-seven-hessian-gap} follows. A local slice
transverse to the compact unitary-group orbit gives the claimed
strict local maximum modulo symmetry.
\end{proof}


\begin{theorem}[quadratic stability near the maximiser; proved from SEH-1--SEH-3, rational constants EB-0--EB-3]
\label{thm:entropy-seven-quantbasin}
Let $\mathcal O_E$ be the local-unitary, global-phase
and qubit-permutation orbit of $\phi=\psi_{x_*}$, the equality set of
Theorem~\ref{thm:seven-entropy-global-maximum}.
For any normalised pure seven-qubit state $z$ write
$\delta(z)=\min_{\phi\in\mathcal O_E}\|z-\phi\|_2$.
If $\delta(z)\le1/20000$, then
\begin{equation}\label{eq:entropy-seven-quantbasin}
\boxed{\qquad
 E_3(z)\le E_3(\psi_{x_*})-\frac{\delta(z)^2}{20}.
 \qquad}
\end{equation}
\end{theorem}
\begin{proof}
Let $\phi'\in\mathcal O_E$ be a closest point to $z$; by symmetry we may take $\phi'=\phi$.
After the allowed global phase its overlap with $z$ is real
and nonnegative. Thus
\[
 z=\cos\theta\,\phi+\sin\theta\,w,\qquad
 \|w\|=1,\quad\langle\phi,w\rangle=0,\quad
 \delta(z)=2\sin(\theta/2).
\]
Here $\theta=0$ is trivial. The nearest-point
condition implies that $w$ is real orthogonal
to every local-unitary and phase tangent vector
at $\phi$. The full-space constrained Hessian
certificate of Theorem~\ref{thm:entropy-seven-local}
then gives, for $e(t)=E_3(\cos t\,\phi+\sin t\,w)$,
\[
 e'(0)=0,\qquad e''(0)<-\frac{239}{1000}.
\]

We obtain a uniform third-derivative bound with no
numerical spectral estimates. For each $|S|=3$,
let $M(t)$ be the $8\times16$ flattening of the
unit vector $\gamma(t)=\cos t\,\phi+\sin t\,w$,
and let $N=M'$, $R=MM^\dagger$. Along this
unit-sphere geodesic,
$M''=-M$, $\|M\|_F=\|N\|_F=1$, so
\[
 \|R'\|_F\le2,\quad \|R''\|_F\le4,\quad
 R'''=-4R',\quad\|R'''\|_F\le8.
\]
At $t=0$ all triple marginals have smallest
eigenvalue at least $1/12$ by the exact
eigenvalue enclosures in SEH-1. As
$\|R'(t)\|_{\rm op}\le2$, Weyl's inequality
gives, for $0\le t\le1/10000$,
\[
 \lambda_{\min}(R(t))
 \ge\frac1{12}-2t
 \ge\frac1{12}-\frac1{5000}
 >\frac1{16}=:m.
\]
The eigenvalues are at most $1$ because $\gamma(t)$
is normalised. Hence
$\|\ln R+I\|_F<6$.

Write
$\mathcal L_R(K)=\int_0^\infty
 (R+sI)^{-1}K(R+sI)^{-1}\,ds$
for the Fr\'echet derivative of $\ln R$.
The resolvent integral bounds imply
\[
 \|\mathcal L_R(K)\|_F\le m^{-1}\|K\|_F,\qquad
 \|D\mathcal L_R[H](K)\|_F
   \le m^{-2}\|H\|_F\|K\|_F.
\]
Differentiating the entropy of one cut three times
(with natural logarithms internally) yields
\[
 \frac{d^3}{dt^3}[-\operatorname{tr}R\log_2 R]
 =-\frac1{\ln2}
 \left[
 \operatorname{tr}R'''(\ln R+I)
 +3\operatorname{tr}R''\mathcal L_R(R')
 +\operatorname{tr}R'\,
       D\mathcal L_R[R'](R')
 \right].
\]
Therefore its absolute value is bounded by
\[
 \frac{8\cdot6+3\cdot4\cdot2\cdot16
                     +2^3\cdot16^2}{\ln2}
 =\frac{2480}{\ln2}<3720,
\]
using $\ln2>2/3$. Averaging the $35$ cuts
gives the same bound $|e'''(t)|<3720$.
Every step holds uniformly along the geodesic.

For $\delta\le1/20000$,
$\theta=2\arcsin(\delta/2)\le2\delta\le1/10000$.
Taylor's theorem and the negative Hessian give
\[
 e(\theta)-e(0)
 \le-\frac{239}{2000}\theta^2
       +\frac{3720}{6}\theta^3
 \le-\left(\frac{239}{2000}-\frac{620}{10000}\right)
       \theta^2
 =-\frac{23}{400}\theta^2
 \le-\frac1{20}\delta^2,
\]
as $\theta\ge\delta$.
EB-0--EB-3 check all needed rational constant
inequalities independently; the spectral and Hessian
inputs are exactly SEH-1--SEH-3. This proves
\eqref{eq:entropy-seven-quantbasin}.
\end{proof}

\begin{corollary}[Global quadratic stability at seven qubits; proved from the certified theorems cited in the proof]\label{cor:seven-global-quadratic}
Let $\delta_p(z)$ be the distance from a normalised pure seven-qubit state $z$ to the orbit of $\Psiico$ under local
unitaries, global phase and qubit permutations, and $\delta(z)$ the distance to $\mathcal O_E$ as in
Theorem~\ref{thm:entropy-seven-quantbasin}. There are constants $c_p,c_e>0$ such that every such $z$ satisfies
\[
 \pi_3(z)-\frac{121}{917}\ge c_p\,\delta_p(z)^2,\qquad
 E_3(\psi_{x_*})-E_3(z)\ge c_e\,\delta(z)^2 .
\]
\end{corollary}
\begin{proof}
Both distances are at most $\sqrt2$, since $\min_\theta\|z-e^{i\theta}\phi\|^2=2-2|\langle z,\phi\rangle|$. If
$\delta_p(z)\le1/250$, a qubit permutation $\sigma$ carries a nearest point of the orbit into the local-unitary and
phase orbit of $\Psiico$; $\pi_3$ is permutation invariant, so Theorem~\ref{thm:icoquantbasin} applied to $\sigma z$
gives the bound with $1/1024$. The set $\{\delta_p\ge1/250\}$ of unit vectors is compact and $\pi_3-121/917$ is
continuous on it. It is positive there, because equality holds only on the orbit
(Corollary~\ref{cor:seven-global-from-fivecopy}), so it has a positive minimum $m_p$. Hence
$c_p=\min\{1/1024,m_p/2\}$ works. The entropy bound follows in the same way from
Theorem~\ref{thm:entropy-seven-quantbasin} for $\delta\le1/20000$, continuity of $E_3$, and the equality case of
Theorem~\ref{thm:seven-entropy-global-maximum}. The constants come from compactness and are not explicit.
\end{proof}

\subsection{Mixed states attain neither optimum}\label{sec:seven-mixed-roofs}
For a seven-qubit density matrix $\rho$ and $|S|=3$ let
\[
 E_S(\rho)=\inf_{\rho=\sum_\ell p_\ell|\psi_\ell\rangle\langle\psi_\ell|}\sum_\ell p_\ell\,S\bigl(\rho_S^{(\ell)}\bigr),\qquad
 L_S(\rho)=\inf_{\rho=\sum_\ell p_\ell|\psi_\ell\rangle\langle\psi_\ell|}\sum_\ell p_\ell\bigl(1-\operatorname{tr}(\rho_S^{(\ell)})^2\bigr)
\]
be the convex roofs of the entropy (in bits) and of the linear entropy of the cut $S|S^c$, where $\rho_S^{(\ell)}$ is the
marginal of $\psi_\ell$.

\begin{lemma}[no flat extremal planes]\label{lem:seven-flat-extremal-planes}
Let $f$ be a real function on the unit sphere of $(\mathbb C^2)^{\otimes7}$, invariant under phases and local unitaries,
and let $\phi$ be a global extremum of $f$ such that $f$ is $C^2$ near $\phi$, every one-qubit marginal of $\phi$ is $I/2$,
and the constrained Hessian $H$ of $f$ at $\phi$ satisfies $H[w,w]\ne0$ for every nonzero real tangent vector
$w$ ($\operatorname{Re}\langle\phi,w\rangle=0$) orthogonal to
\[
 G=\operatorname{span}_{\mathbb R}\{i\phi,\ i\sigma_a^{(j)}\phi\}.
\]
Then no complex subspace of dimension at least two containing $\phi$ has all its unit vectors global extrema of the
same kind.
\end{lemma}
\begin{proof}
Suppose $V\ni\phi$ is such a subspace with $\dim V\ge2$. Pick a unit vector $\psi\in V$ with $\langle\phi,\psi\rangle=0$, and let
$w$ be $\psi$ or $i\psi$. The great circle $t\mapsto\cos t\,\phi+\sin t\,w$ lies in the unit sphere of $V$, so $f$ is
constant on it, and at the critical point $\phi$ its second derivative along this geodesic is $H[w,w]$; hence $H[w,w]=0$.
The vectors spanning $G$ are tangent, as $\langle\phi,\sigma_a^{(j)}\phi\rangle$ is real, and the orbit of $\phi$ consists of
critical points, so differentiating the vanishing constrained gradient along it gives $Hg=0$ for $g\in G$. Writing $w=g+w'$
with $g\in G$ and $w'$ tangent and orthogonal to $G$, $0=H[w,w]=H[w',w']$, so $w'=0$: both $\psi$ and $i\psi$ lie in $G$.

But $G\cap iG=\{0\}$. Let $\mathfrak k$ be the real span of $iI$ and the $i\sigma_a^{(j)}$, and suppose $g=A\phi=iA'\phi$ lies in $G\cap iG$, with
$A,A'\in\mathfrak k$. Put $B=-A'$, so $A\phi=-iB\phi$. Since $A,B$ are anti-Hermitian,
\[
 \langle\phi,[A,B]\phi\rangle=\langle B\phi,A\phi\rangle-\langle A\phi,B\phi\rangle=-2i\|B\phi\|^2 .
\]
The commutator $[A,B]$ is a sum of traceless one-site operators, and every one-qubit marginal of $\phi$ is $I/2$, so the left
side vanishes. Hence $B\phi=0$ and $A\phi=0$. Thus $\psi=0$, a contradiction.
\end{proof}

\begin{proposition}[mixed-state roofs; proved from the certified theorems cited in the proof]\label{prop:seven-mixed-roofs}
Every seven-qubit density matrix $\rho$ satisfies
\[
 \frac1{35}\sum_{|S|=3}E_S(\rho)\le E_3(\psi_{x_*}),\qquad \frac1{35}\sum_{|S|=3}L_S(\rho)\le1-\frac{121}{917}.
\]
Equality holds in the first bound only for pure states local-unitarily and permutation equivalent to $\psi_{x_*}$, and in
the second only for pure states equivalent to $\Psiico$.
\end{proposition}
\begin{proof}
Using one ensemble of $\rho$ for all $35$ infima bounds the left sides by $\sum_\ell p_\ell E_3(\psi_\ell)$ and
$\sum_\ell p_\ell(1-\pi_3(\psi_\ell))$, so Theorem~\ref{thm:seven-entropy-global-maximum} and
Corollary~\ref{cor:seven-global-from-fivecopy} give the bounds and the pure equality cases.

Suppose the entropy bound is attained. The left side is at most the ensemble mean of $E_3(\psi_\ell)$, and each term is
at most $E_3(\psi_{x_*})$, so every member of every ensemble of $\rho$ is a pure maximiser. Every unit
vector $\chi$ in the support of $\rho$ belongs to some ensemble, since $\rho-p|\chi\rangle\langle\chi|\succeq0$ for
$p=\langle\chi,\rho^+\chi\rangle^{-1}$, with $\rho^+$ the inverse of $\rho$ on its support; so every unit vector of the
support is a maximiser. Pick a unit vector $\phi$ in the support; after a local unitary, a qubit permutation and a phase,
applied to $\rho$, which change neither side, $\phi=\psi_{x_*}$. Lemma~\ref{lem:seven-flat-extremal-planes} applies to
$f=E_3$ at $\phi$: near $\phi$ all triple marginals have full rank, so $E_3$ is smooth there, the one-qubit marginals are
$I/2$ (Theorem~\ref{thm:ico}), and Theorem~\ref{thm:entropy-seven-local} gives $H[w,w]<-\frac{239}{1000}\|w\|^2$ for tangent $w$
orthogonal to $G$. So the support has dimension one and $\rho$ is pure.

The linear-entropy case is the same with $f=\pi_3$, a global minimum and $\Psiico$ in place of $\psi_{x_*}$: the Lagrange form
$Q$ of Theorem~\ref{thm:icolocal} is the constrained Hessian, positive semidefinite with kernel exactly $G$.
\end{proof}
Thus mixed-state exclusion follows from two properties of the pure optimiser, maximally mixed one-qubit marginals
and a Hessian nondegenerate transverse to the local-unitary orbit (Lemma~\ref{lem:seven-flat-extremal-planes}); the
four-qubit proof of
Corollary~\ref{cor:fqe-convex-roof} instead uses the square flattenings, which the $3|4$ cuts do not have.
